\chapter{Plan de Pruebas y Certificación}

\section{Procedimientos de prueba}

El aseguramiento de la calidad de la infraestructura de red se ejecutará mediante un protocolo estricto dividido en tres fases: Inspección Visual, Configuración Instrumental y Certificación Electrónica.

\subsection*{Fase 1: Inspección Visual Preliminar}

Antes de conectar el equipo certificador, se realizará una validación física para asegurar que la instalación cumple con las buenas prácticas de manufactura y los estándares ANSI/TIA-568.2-D y ANSI/TIA-569-D.

\begin{itemize}
    \item \textbf{Integridad del Cable:} Verificación de que la chaqueta del cable no presente cortes, abrasiones, quemaduras o deformaciones causadas por una tracción excesiva durante el tendido.
    \item \textbf{Radio de Curvatura:} Confirmación de que ningún mazo de cables exceda el radio de curvatura mínimo (4 veces el diámetro del cable) en las esquinas, subidas a gabinetes o ingresos a cajas de paso.
    \item \textbf{Terminaciones:} Inspección de los Jacks RJ45 en el Patch Panel y en el Área de Trabajo para asegurar que el destrenzado de los pares no supere los 13 mm (0.5 pulgadas) y que el código de colores T568B se haya respetado.
    \item \textbf{Identificación:} Validación de que todos los puntos (Faceplates y Paneles) cuenten con sus etiquetas provisionales o definitivas según la norma ANSI/TIA-606-C.
\end{itemize}

\subsection*{Fase 2: Configuración del Equipo de Certificación}

Se utilizará un analizador de cableado de Nivel IV (ej. Fluke Networks DSX-8000 o equivalente) con calibración vigente.

\begin{itemize}
    \item \textbf{Configuración de Prueba:} Enlace Permanente (Permanent Link).
    \begin{itemize}
        \item \textbf{Justificación:} Se excluyen los patch cords de prueba para medir únicamente la infraestructura fija instalada (del Patch Panel al Jack de usuario), que es la responsabilidad del instalador.
    \end{itemize}
    \item \textbf{Tipo de Cable:} F/UTP Cat 6A LSZH (con blindaje general foil).
    \item \textbf{NVP (Nominal Velocity of Propagation):} Se ajustará el valor NVP según la ficha técnica del fabricante del cable suministrado (típicamente entre 65\% y 72\%) para asegurar una medición precisa de la longitud.
    \item \textbf{Estándar Límite:} TIA Cat 6A Perm. Link (+PoE), para validar no solo datos sino también la capacidad de soportar alimentación eléctrica (PoE) sin sobrecalentamiento.
\end{itemize}

\subsection*{Fase 3: Ejecución de Pruebas de Desempeño (Autotest)}

Se ejecutará la secuencia de Autotest en el 100\% de los \textbf{\totalPuntos{} puntos de red} distribuidos en los \totalPisos{} niveles, evaluando los siguientes parámetros críticos para la operación a 500 MHz:

\begin{enumerate}
    \item \textbf{Mapa de Cableado (Wire Map):} Verifica continuidad pin a pin, cortocircuitos, pares abiertos, pares cruzados y pares divididos (split pairs).
    \item \textbf{Longitud (Length):} Confirma que el enlace no supere los 90 metros.
    \item \textbf{Pérdida de Inserción (Insertion Loss):} Mide la atenuación de la señal a medida que viaja por el cable.
    \item \textbf{Pérdida de Retorno (Return Loss):} Mide la energía reflejada causada por diferencias de impedancia (crítico para 10 Gbps).
    \item \textbf{Diafonía (NEXT - Near End Crosstalk):} Mide la interferencia de un par sobre otro en el extremo cercano.
    \item \textbf{Suma de Potencia de Diafonía (PS NEXT):} Crítico en Cat 6A, mide el efecto combinado de tres pares interfiriendo sobre el cuarto par.
    \item \textbf{ACR-N y ACR-F:} Relación Atenuación-Diafonía en extremos cercano y lejano.
\end{enumerate}

\subsection*{Protocolo de Corrección de Fallas}

En caso de que un punto obtenga el resultado FALLA (FAIL) o PASA (Marginal)*:

\begin{enumerate}
    \item \textbf{Diagnóstico:} Utilizar la función de diagnóstico del equipo (ej. HDTDX/HDTDR) para localizar la distancia exacta de la falla (si es en el conector o a mitad del cable).
    \item \textbf{Acción Correctiva:}
    \begin{itemize}
        \item Falla en extremos: Re-terminar el conector (Jack/Panel) manteniendo el trenzado hasta el punto de contacto.
        \item Falla en cable: Si se detecta daño físico en el trayecto, el cable deberá ser reemplazado en su totalidad (no se aceptan empalmes).
    \end{itemize}
    \item \textbf{Re-certificación:} Una vez corregido, el punto debe ser probado nuevamente y guardado solo cuando el resultado sea ``PASA'' sólido.
\end{enumerate}

\section{Protocolos de certificación}

Para asegurar la conformidad del sistema de cableado estructurado con los requisitos de desempeño para 10GBASE-T, se establecen los siguientes protocolos mandatarios que regirán la aceptación de la obra.

\subsection*{Estándar de Referencia Principal}

El protocolo de prueba se configurará estrictamente bajo los parámetros de la norma ANSI/TIA-568.2-D para la Categoría 6A.

\begin{itemize}
    \item \textbf{Rango de Frecuencia:} Las pruebas de barrido se realizarán desde 1 MHz hasta 500 MHz.
    \item \textbf{Aplicación Soportada:} El protocolo debe validar matemáticamente la capacidad del canal para soportar el estándar IEEE 802.3an (10 Gigabit Ethernet).
\end{itemize}

\subsection*{Topología de Prueba: Enlace Permanente (Permanent Link)}

Se define el Enlace Permanente como la parte fija de la red instalada, excluyendo los patch cords de usuario y de equipos. Este será el único protocolo válido para la certificación de la instalación.

\begin{itemize}
    \item \textbf{Punto de Inicio:} Adaptador de prueba conectado directamente al Patch Panel en el Gabinete (IDF/MDF).
    \item \textbf{Punto Final:} Adaptador de prueba conectado a la toma de telecomunicaciones (Jack RJ45) en el Área de Trabajo.
    \item \textbf{Límite Físico:} La longitud máxima del enlace permanente no debe exceder los 90 metros. Cualquier enlace superior a esta distancia será considerado ``FALLA'' automáticamente por el protocolo.
\end{itemize}

\subsection*{Criterios de Aceptación y Margen de Seguridad}

\begin{itemize}
    \item \textbf{Resultado ``PASA'' (PASS):} Solo se aceptarán enlaces que obtengan este resultado de manera inequívoca.
    \item \textbf{Resultado ``PASA'' (Marginal Pass)*:} Aquellos enlaces cuyo desempeño esté dentro del margen de error del equipo de medición serán considerados condicionales. El protocolo del proyecto exige que estos puntos sean revisados, re-terminados y vueltos a probar para buscar un ``PASA'' sólido, garantizando la longevidad de la instalación.
    \item \textbf{Resultado ``FALLA'' (FAIL):} Implica la reconstrucción obligatoria del enlace (cambio de cable o conectores) sin costo adicional para la entidad.
\end{itemize}

\subsection*{Protocolo de Calibración de Equipos}

El equipo de certificación (Scanner) deberá contar con:

\begin{itemize}
    \item \textbf{Certificado de Calibración:} Emitido por el fabricante o laboratorio autorizado con una antigüedad no mayor a 12 meses.
    \item \textbf{Ajuste de Referencia:} Se deberá establecer una nueva referencia (``Set Reference'') diariamente al inicio de la jornada o cada vez que se cambien los adaptadores de enlace permanente, para eliminar la resistencia interna de los cables de prueba del equipo.
\end{itemize}

\section{Reportes de certificación esperados}

La aceptación final del proyecto estará condicionada a la entrega de un Dossier de Calidad que documente el desempeño de cada uno de los \textbf{\totalPuntos{}} puntos de red instalados en los \totalPisos{} niveles del pabellón. Los reportes deberán cumplir con las siguientes características formales:

\subsection*{Formato de Entrega (Inalterabilidad)}

\begin{itemize}
    \item \textbf{Formato Digital:} Los resultados se entregarán en formato PDF nativo generado directamente por el software de gestión del equipo certificador (ej. LinkWare PC). No se aceptarán resultados exportados a hojas de cálculo (Excel) o formatos de texto editables, para garantizar la integridad y veracidad de los datos.
    \item \textbf{Organización:} El archivo digital estará organizado lógicamente por carpetas correspondientes a cada nivel y gabinete (IDF), facilitando la búsqueda futura de información.
\end{itemize}

\subsection*{Contenido del Reporte Detallado (Por Nodo)}

Se generará una hoja de reporte individual para cada enlace probado, la cual deberá contener visual y numéricamente:

\begin{enumerate}
    \item \textbf{Identificación:} ID del punto (Etiqueta) acorde al estándar ANSI/TIA-606-C (ej. 302-D1).
    \item \textbf{Resumen Pasa/Falla:} Resultado global claro en el encabezado.
    \item \textbf{Gráficas de Frecuencia:} Gráficos detallados de NEXT, Return Loss y Atenuación en el dominio de la frecuencia (1-500 MHz), mostrando la línea límite del estándar y la curva de desempeño medido.
    \item \textbf{Mapa de Cableado:} Representación gráfica de la continuidad de los 8 conductores y la pantalla (si aplica), confirmando la secuencia T568B.
    \item \textbf{Longitud:} Distancia exacta del enlace en metros.
    \item \textbf{Headroom (Margen):} Indicador del margen de seguridad (en dB) por encima del estándar. Un margen alto indica una instalación de excelente calidad.
\end{enumerate}

\subsection*{Información de Trazabilidad y Calibración}

Cada hoja de reporte debe incluir automáticamente los metadatos de la prueba para fines de auditoría:

\begin{itemize}
    \item \textbf{Fecha y Hora:} Momento exacto de la certificación.
    \item \textbf{Instrumento:} Modelo y Número de Serie del equipo certificador utilizado.
    \item \textbf{Calibración:} Fecha de la última calibración del equipo (vigente).
    \item \textbf{Operador:} Nombre del técnico responsable que ejecutó la prueba.
\end{itemize}

\subsection*{Reporte Sumario}

Adicionalmente, se entregará un listado resumen que condense en una tabla los \textbf{\totalPuntos{}} IDs de los puntos, la longitud de cada uno y el estado final (``PASA''), sirviendo como inventario rápido de la infraestructura entregada.

\section{Protocolos de aceptación}

La recepción final de la infraestructura de cableado estructurado y la conformidad del servicio estarán sujetas al cumplimiento estricto del siguiente protocolo de validación, el cual garantiza que la inversión realizada cumple con los estándares para soportar tráfico de 10 Gbps.

\subsection*{Criterios Técnicos de Aceptación (Validación Instrumental)}

La red será aceptada únicamente si el 100\% de los puntos instalados (\textbf{\totalPuntos{} nodos}) cuentan con un certificado de prueba con resultado ``PASA'' (PASS) bajo el estándar TIA-568.2-D Categoría 6A.

Para la aceptación técnica, se revisarán detalladamente los márgenes de desempeño en los siguientes parámetros críticos reportados por el equipo certificador:

\begin{enumerate}
    \item \textbf{NEXT (Near-End Crosstalk):} Se verificará que la diafonía (interferencia) en el extremo cercano se mantenga por encima de la línea límite del estándar a lo largo de todo el barrido de frecuencia (1-500 MHz). Valores marginales o fallas indicarán una mala terminación en el Jack o Patch Panel.
    \item \textbf{PSNEXT (Power Sum NEXT):} Dado que la red operará con transmisiones simultáneas en los 4 pares (necesario para 10 Gigabit Ethernet), la aceptación depende de que la suma de potencias de diafonía no exceda los límites permitidos.
    \item \textbf{Pérdida de Retorno (Return Loss):} Este parámetro es condición sine qua non para la aceptación. Se rechazará cualquier enlace donde la pérdida de retorno sea inferior al estándar, ya que esto provocaría retransmisiones de paquetes y lentitud en la red, a pesar de tener ``buen cable''.
    \item \textbf{Mapa de Cableado y Continuidad:} No se aceptará ningún punto con errores de polaridad, pares invertidos o blindaje discontinuo.
\end{enumerate}

\subsection*{Procedimiento de Verificación en Campo (Muestreo)}

Como parte del acto de recepción de obra, el Supervisor designado por la Facultad podrá solicitar una verificación aleatoria in situ:

\begin{enumerate}
    \item \textbf{Muestra:} Se seleccionará al azar una muestra representativa (10\% de 681 puntos, es decir, 69 nodos) distribuidos entre los \totalPisos{} niveles, con énfasis en el Nivel 3 (mayor densidad de laboratorios).
    \item \textbf{Contra-prueba:} El contratista deberá repetir la prueba de certificación en presencia del Supervisor utilizando el mismo equipo calibrado.
    \item \textbf{Criterio de Rechazo del Lote:} Si durante el muestreo se detecta más de un 5\% de fallas o discrepancias respecto al informe entregado, se rechazará el lote completo y se exigirá una nueva certificación del 100\% de la red a costo del contratista.
\end{enumerate}

\subsection*{Requisitos Documentarios para el Cierre}

La firma del Acta de Recepción Definitiva procederá una vez entregados y aprobados los siguientes documentos:

\begin{itemize}
    \item \textbf{Dossier de Certificación:} Archivo digital nativo y PDF con los \totalPuntos{} reportes individuales ``PASA'' de cobre F/UTP Cat 6A, más los 2 reportes de fibra óptica OM4 del backbone.
    \item \textbf{Planos As-Built (Como construido):} Planos actualizados en AutoCAD/PDF reflejando la ubicación exacta final de cada punto y la ruta real de las canalizaciones, con su codificación correspondiente (según TIA-606-C).
    \item \textbf{Certificado de Garantía:} Documento emitido por el instalador garantizando la mano de obra y los componentes pasivos por un periodo determinado (ej. 12 o 24 meses).
\end{itemize}
